\emph{Conventions.} Units $\hbar = c = 1$. Physical vectors and the Dirac field on $(\hb, \Omega)$ are those of Fact \claimref{synthesis.imported-gtg-dictionary}; the Dirac current of a solution $\Phi$, read in the physical frame at a point, is $J_{\Phi} = \Phi e_{0}\tilde\Phi$, conserved. By the construction of Fact \claimref{synthesis.imported-dimock-dirac-car-algebra-on-a-gravitational-field} the anticommutator pairing of the smeared fields is the conserved inner product of the solutions they define: with $Sf$ the solution of the causal propagator applied to a test spinor $f$, $\{\psi(f), \psi(h)^{*}\} = (Sf, Sh)$, the flux of the current through any Cauchy surface, and $(Sf, \Phi) = \int f^{\dagger}\Phi$ for smooth solutions $\Phi$ (Green's identity). Let $\mathcal{K}$ be the completion of the solutions with compactly supported Cauchy data in this norm; $\psi(\Phi)$, $\Phi \in \mathcal{K}$, is the field of a solution, with $\|\psi(\Phi)\| = \|\Phi\|$ and $\{\psi(\Phi), \psi(\Phi')^{*}\} = (\Phi, \Phi')$. That the Dimock pairing is this inner product, and $(Sf, \Phi) = \int f^{\dagger}\Phi$, is clause (b) of that Fact (Dimock's reproducing formula). $\mathcal{H}$, $\Lambda$, $a$ and the canonical generator field $\sqrt{2a}\,\Lambda\psi$ are those of Theorem \claimref{synthesis.field-strength-cancels-from-the-leading-null-singularity-of-hadamard-dirac-states-on-null-hypersurfaces}, whose frame has $k\cdot e_{0} = 1$ along each generator; for a null direction $k$ with $k\cdot e_{0} \ne 1$ the Hermitian projector is $\Lambda_{k} = e_{0}k/(2\,k\cdot e_{0})$, with flux $2(k\cdot e_{0})\langle(\Lambda_{k}\psi)^{\dagger}\Lambda_{k}\psi\rangle_{0}$. The trace operator below is written $\iota_{\mathcal{H}}$ (not to be confused with the tidal function of that theorem). The field is assumed to carry gauge-invariant quasi-free Hadamard states of the Dirac field near $\mathcal{H}_{c}$, as in that theorem; only (iii) uses them.

\emph{Step 1 (the current is future causal; computed).} Write $\Phi\tilde\Phi = \alpha + I\beta$ (scalar plus pseudoscalar, $I = I_{4}$). Then $Q(J_{\Phi}) = \Phi e_{0}(\alpha + I\beta)e_{0}\tilde\Phi = \Phi(\alpha - I\beta)\tilde\Phi = \alpha^{2} + \beta^{2} \ge 0$, and $J_{\Phi}\cdot e_{0} = \langle\Phi^{\dagger}\Phi\rangle_{0} \ge 0$, with $\Phi^{\dagger} = e_{0}\tilde\Phi e_{0}$; so $J_{\Phi}$ is future causal and $J_{\Phi}\cdot N \ge 0$ for every future causal $N$. Along a null $k$ the flux density is $J_{\Phi}\cdot k = 2(k\cdot e_{0})\langle(\Lambda_{k}\Phi)^{\dagger}\Lambda_{k}\Phi\rangle_{0}$ and depends on $\Lambda_{k}\Phi$ only. The check verifies these identities exactly on sampled spinors, five null directions and four timelike vectors.

\emph{Step 2 (flux inequality; proved).} This is the standard energy inequality for a conserved causal current through a characteristic hypersurface (the energy identity of Fact \claimref{synthesis.imported-haefner-nicolas-dirac-characteristic-problem-on-light-cones-near-the-vertex}(a), \cite{HafnerNicolas2011}, Eq.~(14), and Hörmander's energy method), re-derived here in the store's normalisation. Let $\Phi$ be a smooth solution with compactly supported Cauchy data and $\Sigma$ a Cauchy surface to the future of $\mathcal{H}_{c}$ (a level set of a Cauchy time function above its maximum on the compact $\mathcal{H}_{c}$). Put $\mathcal{R} = J^{+}(\mathcal{H}_{c})\cap J^{-}(\Sigma)$, compact by global hyperbolicity. Its boundary consists of $\Sigma\cap J^{+}(\mathcal{H}_{c})$ and of the achronal, hence Lipschitz, boundary $\partial J^{+}(\mathcal{H}_{c})\cap J^{-}(\Sigma)$, which contains $\mathcal{H}_{c}$ because $\mathcal{H}\cap N$ is achronal and $N$ is causally convex. The divergence theorem for the conserved current gives: the flux of $J_{\Phi}$ out through $\Sigma\cap J^{+}(\mathcal{H}_{c})$ equals the flux in through $\partial J^{+}(\mathcal{H}_{c})$. By Step 1 every piece of the second flux is non-negative, and so is the flux through the rest of $\Sigma$. The flux measure on $\mathcal{H}$ is $(J_{\Phi}\cdot k)\,a\,du\,d^{2}y$. Hence
\[
\int_{\mathcal{H}_{c}} 2a\,|\Lambda\Phi|^{2}\,du\,d^{2}y \;\le\; \int_{\Sigma\cap J^{+}(\mathcal{H}_{c})} J_{\Phi}\cdot d\Sigma \;\le\; \|\Phi\|^{2}.
\]
No characteristic Cauchy problem is used: only conservation, causality of the current, and the divergence theorem on a region with Lipschitz boundary.

\emph{Step 3 (the contraction, the field, every state; proved).} By Step 2 the trace $\iota_{\mathcal{H}}^{*}\Phi = \sqrt{2a}\,\Lambda\Phi|_{\mathcal{H}_{c}}$, defined on the dense set of smooth solutions with compactly supported Cauchy data, satisfies $\|\iota_{\mathcal{H}}^{*}\Phi\| \le \|\Phi\|$ and extends uniquely to a contraction $\iota_{\mathcal{H}}^{*}: \mathcal{K} \to L^{2}(\mathcal{H}_{c})$; its adjoint $\iota_{\mathcal{H}}$ is a contraction. For smooth $F$ and smooth $\Phi$, $(\iota_{\mathcal{H}}F, \Phi) = \langle F, \iota_{\mathcal{H}}^{*}\Phi\rangle = \int F^{\dagger}\sqrt{2a}\,\Lambda\Phi\,du\,d^{2}y$, the generator field smeared with $F$ paired with $\Phi$; so the generator field smeared with $F$ is the field of the solution $\iota_{\mathcal{H}}F$, $\psi_{\mathcal{H}}(F) := \psi(\iota_{\mathcal{H}}F)$. This is the characteristic-data embedding $\psi_{\mathcal{H}}(F) = \psi(\Phi_{F})$, with $\Phi_{F} = \iota_{\mathcal{H}}F$. It is bounded, $\|\psi_{\mathcal{H}}(F)\| = \|\iota_{\mathcal{H}}F\| \le \|F\|$, and $\{\psi_{\mathcal{H}}(F), \psi_{\mathcal{H}}(G)^{*}\} = (\iota_{\mathcal{H}}F, \iota_{\mathcal{H}}G) = \langle F, \iota_{\mathcal{H}}^{*}\iota_{\mathcal{H}}G\rangle$ with $0 \le \iota_{\mathcal{H}}^{*}\iota_{\mathcal{H}} \le 1$. For any state $\omega$ and $A = \psi_{\mathcal{H}}(F)$, $0 \le \omega(AA^{*}) \le \omega(AA^{*} + A^{*}A) = \|\iota_{\mathcal{H}}F\|^{2}$, and the same for $\omega(A^{*}A)$; by Cauchy-Schwarz $|\omega(\psi_{\mathcal{H}}(F)\psi_{\mathcal{H}}(G)^{*})| \le \|F\|\,\|G\|$. So the two restricted two-point functions are bounded positive operators $W^{\pm}_{\mathcal{H},\omega}$ on $L^{2}(\mathcal{H}_{c})$, and exactly $W^{+}_{\mathcal{H},\omega} + W^{-}_{\mathcal{H},\omega} = \iota_{\mathcal{H}}^{*}\iota_{\mathcal{H}} \le 1$ (the anticommutator is state independent). This is (i).

\emph{Step 4 (generator scaling; proved).} Fix transverse profiles $f, g$ and write $w(u, u')$ for the $(f, g)$ component of $W^{\pm}_{\mathcal{H},\omega}$, an operator on $L^{2}(du)$ of norm at most $\|f\|\,\|g\|$ by Step 3. For test functions $\varphi, \chi$ in $x$ put $(U_{\lambda}\varphi)(u) = \lambda^{-1/2}\varphi((u - u_{0})/\lambda)$, unitary. The substitution $u = u_{0} + \lambda x$, $u' = u_{0} + \lambda x'$ gives $\lambda\int\!\!\int\bar\varphi(x)\,w(u_{0} + \lambda x, u_{0} + \lambda x')\,\chi(x')\,dx\,dx' = \langle U_{\lambda}\varphi, w\,U_{\lambda}\chi\rangle$ (checked exactly for polynomial kernels and test functions, with the weight $\tfrac12$ per leg and not $1$). So the scaled kernel is $U_{\lambda}^{*}wU_{\lambda}$, of the same norm, for every $\lambda$ for which the supports stay in $\mathcal{H}_{c}$: $\lambda^{\varpi}w(u_{0} + \lambda\cdot, u_{0} + \lambda\cdot) \to 0$ for $\varpi > 1$, and by weak-operator compactness of the unit ball the scaling limit points exist and are positive contractions. Which of them is the null vacuum is the question of the generator scaling limit of Theorem \claimref{synthesis.microlocal-condition-on-a-generator-gives-the-local-null-ground-state}(ii); the bound does not settle it. This is (ii). The same bound holds for $W_{0}\otimes\id_{K}$, the dilation-invariant positive-frequency projection, and smooth pullbacks vanish at rate $\lambda$. So any restricted parametrix that agrees with $W_{\mathcal{H},\omega}$ up to a smooth pullback has $R^{\mathrm{par}}$ of generator scaling degree at most one.

\emph{Step 5a (Hadamard states differ by smooth pullbacks; proved, no (Iso)).} Let $\{\mathcal{H}_{v}\}$ be the foliation of (iii) (for example generated from a smooth family of spacelike cuts displaced along $n$), with coordinates $(u, y, v)$ and contractions $\iota_{v}$ as in Step 3. The thin-layer spinor $f_{n} = \int\rho_{n}(v)\,\delta_{\mathcal{H}_{v}}\otimes F_{v}\,dv$, where $\langle\delta_{\mathcal{H}_{v}}\otimes F_{v}, \Phi\rangle = \langle F_{v}, \iota_{v}^{*}\Phi\rangle$, is smooth with compact support: it is $\rho_{n}(v)$ times a smooth spinor field. For smooth $\Phi$, $(Sf_{n}, \Phi) = \int\rho_{n}(v)\,(\iota_{v}F_{v}, \Phi)\,dv$, so $Sf_{n} = \int\rho_{n}(v)\,\iota_{v}F_{v}\,dv$ as a Bochner integral in $\mathcal{K}$ (the integrand is bounded by $\sup_{v}\|F_{v}\|$ and weakly measurable). The map $v \mapsto \iota_{v}F_{v}$ is weakly continuous: $(\iota_{v}F_{v}, \Phi) = \int F_{v}^{\dagger}\sqrt{2a_{v}}\,\Lambda_{v}\Phi|_{\mathcal{H}_{v}}\,du\,d^{2}y$ is continuous in $v$ for the dense set of smooth $\Phi$, and the norms are bounded. Hence $Sf_{n} \to \iota_{\mathcal{H}}F$ weakly, with no further hypothesis. Now let $\omega_{1}, \omega_{2}$ be gauge-invariant quasi-free Hadamard states. By Fact \claimref{synthesis.imported-dirac-hadamard-form-propagates-and-adiabatic-states-are-hadamard}(a) their kernels near $\mathcal{H}_{c}$ are $G_{\mathrm{par}} + D^{\omega_{i}}$ with $D^{\omega_{i}}$ smooth, $G_{\mathrm{par}}$ the Dirac operator applied to a Hadamard parametrix of its square (Fact \claimref{synthesis.imported-hadamard-parametrix-of-the-squared-dirac-operator-leads-with-van-vleck-root-times-spin-transport}); so the difference $b = W^{\pm}_{\omega_{1}} - W^{\pm}_{\omega_{2}}$ is a smooth kernel $D = D^{\omega_{1}} - D^{\omega_{2}}$ near $\mathcal{H}_{c}$ (smooth differences: Fact \claimref{synthesis.imported-microlocal-hadamard-condition}(b)). As a sesquilinear form on $\mathcal{K}$, $b$ is bounded (each two-point form is bounded by $1$). For each fixed $n$, $b(Sf_{n}, \cdot)$ is a bounded functional, so $b(Sf_{n}, Sg_{m}) \to b(Sf_{n}, \iota_{\mathcal{H}}G)$ as $m \to \infty$; then $b(\cdot, \iota_{\mathcal{H}}G)$ is bounded, so $b(Sf_{n}, \iota_{\mathcal{H}}G) \to b(\iota_{\mathcal{H}}F, \iota_{\mathcal{H}}G)$. On the other side, $b(Sf_{n}, Sg_{m}) = D(f_{n}, g_{m}) = \int\!\!\int\rho_{n}(v)\rho_{m}(v')\,\langle F_{v}\otimes G_{v'}, D_{v,v'}\rangle\,dv\,dv'$, with $D_{v,v'}$ the compressed, canonically normalised restriction of $D$ to $\mathcal{H}_{v}\times\mathcal{H}_{v'}$, continuous in $(v, v')$; its iterated limit is $\langle F\otimes G, D_{0,0}\rangle$. Therefore $W^{\pm}_{\mathcal{H},\omega_{1}} - W^{\pm}_{\mathcal{H},\omega_{2}} = D_{0,0}$, a genuine pullback of a smooth kernel. Consequently $W^{\mathrm{par}}_{\mathcal{H}} := W_{\mathcal{H},\omega} - D^{\omega}_{0,0}$ is the same for every gauge-invariant quasi-free Hadamard state $\omega$. This is (iii-a). It makes (iii) of Theorem \claimref{synthesis.field-strength-cancels-from-the-leading-null-singularity-of-hadamard-dirac-states-on-null-hypersurfaces} unconditional: the characteristic-data embedding that result cites from memory is Step 3, and the smooth-pullback statement is this step.

\emph{Step 5b (thin layers under (Iso); proved given (Iso)).} Under (Iso), $\|\iota_{v}F_{v}\| = \|F_{v}\|$ is continuous in $v$, so $\|\iota_{v}F_{v} - \iota_{\mathcal{H}}F\|^{2} = \|F_{v}\|^{2} + \|F\|^{2} - 2\,\mathrm{Re}\,(\iota_{v}F_{v}, \iota_{\mathcal{H}}F) \to 0$: the map is norm continuous at $v = 0$. Hence $\|Sf_{n} - \iota_{\mathcal{H}}F\| \le \int\rho_{n}(v)\|\iota_{v}F_{v} - \iota_{\mathcal{H}}F\|\,dv \to 0$ and $\|\psi(f_{n}) - \psi_{\mathcal{H}}(F)\| \to 0$. States are norm continuous, so $W^{\pm}_{\omega}(f_{n}, g_{n}) \to W^{\pm}_{\mathcal{H},\omega}(F, G)$ for every state: the restriction defined by the embedding is the joint thin-layer limit of the bulk two-point functions. Without (Iso) the convergence of Step 5a is weak, and $\|Sf_{n}\| \to \|\iota_{\mathcal{H}}F\|$ is exactly what decides norm convergence. For a Hadamard state, $D^{\omega}(f_{n}, g_{n}) \to \langle F\otimes G, D^{\omega}_{0,0}\rangle$ jointly (continuity of $D^{\omega}_{v,v'}$), so $\lim_{n}G_{\mathrm{par}}(f_{n}, g_{n})$ exists and equals $W^{\mathrm{par}}_{\mathcal{H}}(F, G)$ of Step 5a: under (Iso) the restricted parametrix has this regularisation-free reading, independent of the foliation and the layers. Also under (Iso) $W^{+}_{\mathcal{H},\omega} + W^{-}_{\mathcal{H},\omega} = 1$, so, with $\delta - W_{0}$ the other ordering of the null vacuum, the sum of the two orderings of $R^{\mathrm{par}}$ is $-(D^{\omega,+}_{0,0} + D^{\omega,-}_{0,0})$, a smooth pullback. This is (iii-b).

\emph{Step 6 (where (Iso) is known).} (a) Light cones. For the Dirac field on an arbitrary spacetime and the future light cone $\mathcal{C}^{+}_{0}$ of a point $p_{0}$, within $D_{T} = J^{+}(p_{0})\cap\{0 \le t \le T\}$ in a geodesically convex, globally hyperbolic neighbourhood of $p_{0}$, \cite{HafnerNicolas2011} prove (Fact \claimref{synthesis.imported-haefner-nicolas-dirac-characteristic-problem-on-light-cones-near-the-vertex}) that the trace operator from data on $\Sigma_{T}$ to the transverse components on the cone is an isometry for the $L^{2}$ norms (their Theorem 1, from the energy identity of the conserved current, their Eq.~(14)) and an isomorphism of $H^{1}$ onto the space of characteristic data (their Theorem 2). An isometry with dense range is unitary, so the data-to-solution map is an isometry, $\iota^{*}\iota = 1$, for data on the part of the cone below $t = T$; the two-component transverse data are, up to the frame, $\Lambda\Phi$. A foliation by the cones of the points of a timelike curve satisfies (Iso) leaf by leaf, on the part of each cone covered by the theorem. (b) Flat null planes, any mass. In the conventions of Theorem \claimref{synthesis.bulk-hadamard-states-restrict-to-the-generator-condition-on-flat-null-planes}, whose (i) is the generator CAR on the plane, (Iso) holds. For the thin-layer reading in momentum space, $\langle f_{n}, Sg_{n}\rangle = \tfrac12\int\hat F^{\dagger}\hat G\,|\hat\rho_{n}(p\cdot n_{v})|^{2}\,dp_{+}\,d^{2}p_{\perp}$ over both shells, $n_{v} = \tfrac12(e_{0} - e_{3})$, because the compressed numerator is $\Lambda(p \pm m_{0})e_{0}\Lambda = (p\cdot k)\Lambda$ and $d^{3}p/(2|p_{0}|) = dp_{+}\,d^{2}p_{\perp}/(2|p_{+}|)$. Only $|\hat\rho_{n}| \le 1$ and $\hat\rho_{n} \to 1$ pointwise are used, which hold for every positive approximate identity $\rho_{n}$; dominated convergence then gives $\langle F, G\rangle$. The check samples the defect $1 - \hat\rho_{n}^{2}$ for the transform $n^{2}/(n^{2} + q^{2})$ of the two-sided exponential, an example of these properties (not a smooth layer). (c) Other null hypersurfaces, in particular local Rindler horizons, whose expansion and shear vanish at the point, so that they are not light cones. (Iso) is the solvability, with norm equal to the data norm, of the characteristic Cauchy problem for the Dirac field with data on a general null hypersurface. A route: a compact piece $\mathcal{H}_{c}$ of an achronal null hypersurface lies in an achronal Lipschitz, weakly spacelike Cauchy hypersurface (for example the boundary of $J^{+}(\mathcal{H}_{c})\cup J^{+}(\Sigma_{-})$ for a Cauchy surface $\Sigma_{-}$ to its past), and \cite{HafnerNicolas2011}, Sec.~4.5, treat such hypersurfaces (their weak spacelike condition is non-strict, so null portions are allowed) on spatially compact spacetimes, by Hörmander's energy method, with the energy identity for smooth solutions (their Prop.~4.13). Their Theorem 3 there was not read against the text, and the localisation to a spatially compact field is not carried out (cited from memory), so (Iso) is a condition of (iii-b) there.

\emph{Step 7 (relation to the regularised restriction; argued).} (HS) and (ii) of Theorem \claimref{synthesis.field-strength-cancels-from-the-leading-null-singularity-of-hadamard-dirac-states-on-null-hypersurfaces} define the restriction of the parametrix by regularising $\sigma$ with $i\epsilon s$, $s$ from a time function, at points of $\mathcal{H}$; Step 5a defines $W^{\mathrm{par}}_{\mathcal{H}}$ through a Hadamard state, and Step 5b reads it as a thin-layer limit. Expand the projected kernel in the label difference $\eta$ along $\mathcal{H}$. (1) Exact input: for the Hermitian projector of any null $k$, $\Lambda p\,e_{0}\Lambda\,(k\cdot e_{0}) = (p\cdot k)\Lambda$ for every vector $p$, which at $k\cdot e_{0} = 1$ is the identity of Theorem \claimref{synthesis.bulk-hadamard-states-restrict-to-the-generator-condition-on-flat-null-planes}(ii) (the check takes five regulators and the Hermitian partner $n$ for five null directions). For a complex time-function regulator the numerator shift divided by the shift of $\sigma$ therefore gives a coefficient of $(u - u' + i0)^{-1}\delta^{2}(y - y')$ in which $p\cdot k$ cancels. (2) Argued: in the thin-layer reading the displacement along $n$ is real, $\sigma \approx (u - u')(v - v') - \tfrac12 Q(w)$, and the $i0$ comes from the time function of $G_{\mathrm{par}}$ itself, $s$ depending on both $u - u'$ and $v - v'$; the sign of the $i0$, and with it whether the term is $(u - u' + i0)^{-1}$, depends on $\mathrm{sgn}(v - v')$, and that the layer average reproduces the same coefficient is not shown. (3) Argued: the degree $-2$ part in $\eta$ gives $\log\epsilon$ in one reading and $\log|v - v'|$ in the other, with the same coefficient (the mean of that part over the ellipse of $M$); given (Iso) the thin-layer limit exists by Step 5b, so this coefficient would vanish and the $\log\epsilon$ terms cancel. (4) Parts of degree above $-2$ are integrable at $\eta = 0$ and give the same limit in both readings by dominated convergence; the two definitions may then differ only by local terms $c(u, u')\,\delta^{2}(y - y')$ with $c$ bounded near the diagonal, of generator scaling degree at most $0$, and (HS) would have the same truth value in either. What is exact is $W^{+}_{\mathcal{H},\omega} + W^{-}_{\mathcal{H},\omega} = \iota_{\mathcal{H}}^{*}\iota_{\mathcal{H}}$ (Step 3), so under (Iso) the sum of the two orderings of $R^{\mathrm{par}}$ is a smooth pullback (Step 5b); the log terms of each ordering separately are not controlled by (Iso). This is (v).

\emph{Literature (context; not used in the proofs of (i) to (iii)).} Kay and Wald \cite{KayWald1991} restrict, for the scalar field on spacetimes with a bifurcate Killing horizon, the two-point function of a Hadamard state invariant under the Killing flow to the horizon. The restriction is taken through solutions with characteristic data on the horizon, and the result is the form $-\pi^{-1}\lim_{\epsilon\to 0}\int F_{1}F_{2}(U_{1} - U_{2} - i\epsilon)^{-2}$ in the affine parameter, for data of compact support. Here this is read through the account of \cite{DappiaggiMorettiPinamonti2009}, Sec.~3.2, not against the original. As noted by \cite{Sanders2013}, Sec.~1, their argument rests on the characteristic (Goursat) problem, whose solutions may fail to be smooth across the horizon and for which only local results on special null hypersurfaces were available. Their argument uses the Killing invariance and the global horizon, and is not a local restriction theorem for general null hypersurfaces or for the Dirac field. Gérard and Wrochna \cite{GerardWrochna2016} go the other way, for the scalar field in the interior of a light cone: boundary states on the cone, with a microlocal condition, induce Hadamard bulk states. The bulk-to-boundary maps of Dappiaggi, Moretti and Pinamonti are for scalar fields. For the Dirac field the characteristic problem on a light cone is \cite{HafnerNicolas2011}, used in Step 6.

\emph{Status.} (i), (ii) and (iii-a) are proved, with no characteristic Cauchy problem; (iii-a) assumes that Hadamard states exist. (iii-b) is proved given (Iso). (iv) imports (Iso) on light cones from Fact \claimref{synthesis.imported-haefner-nicolas-dirac-characteristic-problem-on-light-cones-near-the-vertex} (\cite{HafnerNicolas2011}, Theorems 1 and 2), recalls the stored generator CAR on flat null planes, and cites the general null case from memory. \emph{Note (added 2 November).} On every small piece of a geodesic null hypersurface, local Rindler horizons included, (Iso) is (ii) of Theorem \claimref{synthesis.dirac-causal-propagator-between-nearby-null-leaves-tends-to-the-generator-car-on-every-geodesic-null-hypersurface}, derived from the causal propagator and (i) here; it is cited in this note only, not as a premise, since that theorem uses this result. With its (iii)(b) and Theorem \claimref{synthesis.restricted-dirac-hadamard-parametrix-on-a-null-hypersurface-is-the-null-vacuum-plus-a-bounded-kernel} the comparison of Step 7 is settled up to a bounded kernel containing a smooth time-function-dependent term. (v) is argued, apart from the projector identity and $W^{+} + W^{-} = \iota_{\mathcal{H}}^{*}\iota_{\mathcal{H}}$, which are exact. Transport-absent.
